\documentclass[10pt]{article}
\usepackage[T1]{fontenc}
\usepackage[utf8]{inputenc}
\usepackage[french]{babel}
\usepackage{amsmath,amssymb,mathrsfs,mathpazo}
\usepackage{geometry}
\geometry{a4paper,margin=2cm}
\usepackage{graphicx,array,tabularx,tabularray,colortbl,multirow,multicol,booktabs,enumitem,pifont,ulem,fancybox,diagbox,cancel}
\usepackage{pstricks,pstricks-add,pst-plot,pst-tree,pst-node,pst-text,pst-func,pst-math,pst-bspline,pst-3dplot,pst-eucl,pst-solides3d}
\let\pscalculate\undefined
\usepackage{pst-bezier}
\usepackage{tikz,pgfplots,tkz-tab}
\usetikzlibrary{arrows,arrows.meta,patterns,calc,positioning,trees,decorations.markings,shapes.geometric,shapes.misc,shapes.arrows,intersections,decorations.pathmorphing,decorations.pathreplacing,backgrounds,shadows,babel}
\pgfplotsset{compat=1.18}
\usepgfplotslibrary{fillbetween}
\usepackage[np]{numprint}
\usepackage{listings,makecell,pas-tableur,twoopt}
\let\barre\undefined
\usepackage[pstricks]{tablvar}
\definecolor{darkgreen}{RGB}{0,100,0}
\providecommand{\red}{\color{red}}
\providecommand{\blue}{\color{blue}}
\providecommand{\green}{\color{green}}
\providecommand{\black}{\color{black}}
\providecommand{\white}{\color{white}}
\providecommand{\gray}{\color{gray}}
\providecommand{\cyan}{\color{cyan}}
\providecommand{\magenta}{\color{magenta}}
\providecommand{\yellow}{\color{yellow}}
\providecommand{\violet}{\color{violet}}
\providecommand{\orange}{\color{orange}}
\providecommand{\euro}{\texteuro}
\providecommand{\EUR}{\texteuro}
\providecommand{\ieme}{\textsuperscript{e}}
\providecommand{\degres}{\ensuremath{^\circ}}
\providecommand{\degree}{\ensuremath{^\circ}}
\providecommand{\vv}[1]{\overrightarrow{#1}}
\providecommand{\Coord}[1]{#1}
\providecommand{\decosix}{}
\providecommand{\starredbullet}{\ensuremath{\bullet}}
\pagestyle{empty}
\setlength{\parindent}{0pt}
\usepackage{adjustbox}\DeclareUnicodeCharacter{2236}{:}\DeclareUnicodeCharacter{2190}{\ensuremath{\leftarrow}}
\providecommand{\R}{\mathbb{R}}
\providecommand{\N}{\mathbb{N}}
\providecommand{\Z}{\mathbb{Z}}
\providecommand{\Q}{\mathbb{Q}}
\providecommand{\D}{\mathbb{D}}
\providecommand{\C}{\mathbb{C}}
\providecommand{\e}{{\mathrm{e}}}
\providecommand{\vect}[1]{\overrightarrow{#1}}
\providecommand{\vectt}[1]{\overrightarrow{\mathrm{#1}}}
\providecommand{\vvt}[1]{\overrightarrow{\mathrm{#1}}}
\providecommand{\np}[2][]{{#2~#1}}
\providecommand{\vv}[1]{{\overrightarrow{#1}}}
\renewcommand{\d}{{\mathrm{d}}}
\providecommand{\degres}{{^\circ}}
\providecommand{\degre}{{^\circ}}
\providecommand{\blue}{\color{blue}}
\providecommand{\red}{\color{red}}
\providecommand{\enskip}{\quad}
\providecommand{\llbracket}{\lbrack\!\lbrack}
\providecommand{\rrbracket}{\rbrack\!\rbrack}
\providecommand{\euro}{{\text{€}}}
\providecommand{\no}{{\text{n°}}}
\providecommand{\interval}[2]{{[#1;#2]}}
\providecommand{\intervFO}[2]{{[#1;#2[}}
\providecommand{\dotfill}{\ldots}

\providecommand{\ds}{\displaystyle}
\providecommand{\textding}[1]{#1}
\providecommand{\Oij}{\(\mathrm{O};\vec{\imath},\vec{\jmath}\)}
\providecommand{\Oijk}{\(\mathrm{O};\vec{\imath},\vec{\jmath},\vec{k}\)}
\providecommand{\Ouv}{\(\mathrm{O};\vec{u},\vec{v}\)}
\providecommand{\intervFF}[2]{\ensuremath{[#1;#2]}}
\renewcommand{\euro}{\ifmmode\text{\texteuro}\else\texteuro\fi}
\DeclareUnicodeCharacter{00B0}{\ifmmode^{\circ}\else\textdegree\fi}
\DeclareUnicodeCharacter{20AC}{\ifmmode\text{\texteuro}\else\texteuro\fi}
% Énoncés : Éducation nationale. Transcription et corrigés : APMEP, Denis Vergès et contributeurs. Découpage et classement : CoopMaths, Sébastien Lozano, Éric Elter et contributeurs. Adaptation de la présentation : Maths Explorer.
% CC BY-SA : https://coopmaths.fr/www/about/
% https://coopmaths.fr/alea/static/bac/2023/tex/bac_2023_08_sujet1_caledonie_2_cor.tex
\begin{document}
\noindent\textit{APMEP / CoopMaths. CC BY-SA. Présentation adaptée par Maths Explorer.}\par\medskip
\begin{enumerate}
	\item
	\begin{enumerate}
		\item \[u_1 = 5u_0 - 4 \times 0 - 3 = 5 \times 3 - 4 - 3 = 15 - 7 = \boxed{12}\]
		\item \[u_2 = 5u_1 - 4 \times 1 - 3 = 5 \times 12 - 4 - 3 = 60 - 7 = \boxed{53}\]
		\item Il semble que la suite $\left(u_n\right)$ soit croissante et tende vers $+\infty$.
	\end{enumerate}
	\item
	\begin{enumerate}
		\item Soit $P_n$ la proposition $u_n \geqslant n+1$.
		
\textbf{Initialisation} : $u_0 = 3$ et $0+1 = 1$.
		
$3 \geqslant 1$. La proposition est donc vraie au rang $n=0$.\\
\smallskip
		
\textbf{Hérédité} : on suppose la proposition vraie au rang $n \in \mathbb{N}$,\: $u_n \geqslant n + 1$ (hypothèse de récurrence). On va vérifier qu'alors elle est vraie au rang suivant.
\begin{eqnarray*}
u_n \geqslant n + 1 &\iff& 5u_n \geqslant 5(n+1)\\
&\iff& 5u_n - 4n - 3 \geqslant 5n + 5 - 4n - 3\\
&\iff& u_{n+1} \geqslant n + 2 = (n+1) + 1
\end{eqnarray*}
La proposition est donc héréditaire.\\
		\smallskip
		
\emph{Conclusion} : la proposition $P_n$ est vérifiée au rang $n = 0$ et est héréditaire, donc, d'après le principe de récurrence, elle est vraie pour tout entier naturel $n$ : $u_n \geqslant n + 1$.
		\item On a : $\lim\limits_{n \to +\infty}(n+1) = +\infty$. Puisque $u_n \geqslant n+1$, par comparaison, on a :
		\[\boxed{\lim_{n \to +\infty}u_n = +\infty}\]
	\end{enumerate}
	\item
    \begin{enumerate}
		\item\
		\begin{eqnarray*}
			v_{n+1} &=& u_{n+1}-(n+1) - 1\\
			&=& 5u_n - 4n - 3 - n - 1 - 1\\
			&=& 5u_n - 5n - 5\\
			&=& 5\left(u_n - n - 1\right)\\
			&=& 5v_n
		\end{eqnarray*}
La suite $\left(v_n\right)$ est donc une suite géométrique de raison $q=5$ et de premier terme $v_0 = u_0 - 0 - 1 = 2$.
		\item Puisque $\left(v_n\right)$ est une suite géométrique de raison $5$ et de premier terme $v_0=2$, on a :
		\[\boxed{v_n = 2 \times 5^n}\]
		\item $v_n = u_n - n - 1 \iff u_n = v_n + n + 1$. Donc :
		\[\boxed{u_n = 2 \times 5^n + n + 1}\]
		\item  Puisque $5 > 1$, la suite de terme général $5^n$ est strictement croissante, donc $5^{n+1} \geqslant 5^n$.
		\begin{eqnarray*}
			5^{n+1} \geqslant 5^n &\iff& 2 \times 5^{n+1} \geqslant 2 \times 5^n\\
			&\iff& 2\times 5^{n+1} + (n+1) + 1 \geqslant 2 \times 5^n + (n + 1) + 1\\
			&\iff& u_{n+1} \geqslant 2 \times 5^n + n + 2\\
			&\iff& u_{n+1} \geqslant 2 \times 5^n + n + 1\\
			&\iff& u_{n+1} \geqslant u_n
		\end{eqnarray*}
		La suite $\left(u_n\right)$ est donc croissante.
	\end{enumerate}
	\item
	\begin{enumerate}
		\item\
		
		\fbox{\parbox{.333\textwidth}{
		\begin{ttfamily}
		def suite():\par
		\hspace{0.5cm}u=3\par
		\hspace{0.5cm}n=0\par
		\hspace{0.5cm}while \textbf{u<10**7}:\par
		\hspace{1cm}u= \textbf{5*u-4*n-3}\par
		\hspace{1cm}n=n+1\par
		\hspace{0.5cm}return n
		\end{ttfamily}}}
		\item\
		\begin{center}
		\begin{tabularx}{.5\textwidth}{|*{3}{>{\centering\arraybackslash}X|}}
			\hline
			\tt{u} & \tt{n} & \tt{u < 10\up{7}} \\ \hline
			3 & 0 & VRAI \\ \hline
			12 & 1 & VRAI \\ \hline
			53 & 2 & VRAI \\ \hline
			254 & 3 & VRAI \\ \hline
			\np{1255} & 4 & VRAI \\ \hline
			\np{6256} & 5 & VRAI \\ \hline
			\np{31257} & 6 & VRAI \\ \hline
			\np{156258} & 7 & VRAI \\ \hline
			\np{781259} & 8 & VRAI \\ \hline
			\np{3906260} & 9 & VRAI \\ \hline
			\np{19531261} & 10 & FAUX \\ \hline
		\end{tabularx}
		\end{center}
La valeur renvoyée par cette fonction est $n = 10$. C'est le rang à partir duquel $u_n \geqslant 10^{7}$.
	\end{enumerate}
\end{enumerate}

\bigskip

\end{document}
